\section{Conclusion}
\label{sec:conclusion}
This study evaluates LLM-derived traveler models through finite responses rather than baseline imitation alone. On the controlled benchmark, response-aware supervision improves Teacher fidelity and a logit choice model with independent neural timing performs better still on response error. That ranking does not carry over to the stated-choice samples: the Teacher-best Student reverses a human accessibility response. In the controlled Helsinki execution experiment, improved routability likewise does not ensure closer transmission of the predicted response.

The practical contribution is a response-centric validation procedure linking compression, empirical comparison and network execution. Each level answers a different question, and success at one cannot stand in for evidence at another. A traveler model intended for a specified intervention should therefore be evaluated on its response at all relevant levels, with baseline demand and computational cost reported separately. The current evidence establishes these distinctions within the tested Teacher, samples and execution setting; broader forecasting claims require the corresponding out-of-setting tests.
